\section*{Заключение}

В ходе выполнения работы были получены следующие результаты.

\begin{itemize}
    \item Проведён обзор существующих инструментов для моделирования и анализа мультиагентных систем, включая MATLAB Sensor Fusion and Tracking Toolbox, Repast (Agent-Based Modeling Toolkit) и SIMLAN, с точки зрения их применимости к задачам трекинга.
    
    \item Собраны и систематизированы требования от членов научной группы, разрабатывающих современные мультиагентные алгоритмы трекинга. Требования структурированы по ключевым группам: генерация сценариев (сенсоры, объекты, помехи), поддержка алгоритмов (базовые реализации и возможность интеграции пользовательских), а также средства сравнительного анализа (визуализация ошибок и траекторий).
    
    \item Спроектирована модульная архитектура системы, обеспечивающая разделение компонентов симуляции, выполнения алгоритмов и визуализации. Разработан пользовательский сценарий, в рамках которого исследователь задаёт параметры эксперимента, запускает вычислительный конвейер и получает наглядное сравнение эффективности алгоритмов.
    
    \item Реализована программная система в виде веб-приложения на основе фреймворка Django, обеспечивающая доступ к функциональности через удобный пользовательский интерфейс.
    
    \item Интегрированы эталонные мультиагентные алгоритмы трекинга: алгоритм на основе SPSA, его ускоренная версия и распределённый фильтр Калмана, используемые в качестве базовой линии для сравнения новых алгоритмов.
    
    \item Результаты работы представлены на XXVII конференции молодых учёных «Навигация и управление движением», а также опубликованы в журнале IEEE Transactions on Automatic Control.
\end{itemize}

\noindent Исходный код разработанной системы: \textit{[добавьте ссылку на репозиторий]}.